\begin{textsample}{POS Dim 3 – 2009\_c4 – Score 1.00 – 2009\_c4/enade\_2009\_ci\_econ\_d1cc84d7e1.txt}  \label{ex:f3_pos_125}
Capítulo 4 Percepção sobre a Prova As percepções dos estudantes sobre a prova aplicada no ENADE/2009 na área de CIÊNCIAS ECONÔMICAS foram mensuradas por meio de nove itens que avaliaram desde o grau de dificuldade da prova até o tempo gasto para concluí -la.
As questões foram analisadas separando -se concluintes e ingressantes, e as percepções sobre a prova foram relacionadas com o desempenho dos estudantes e com a Região Geográfica de localização do curso6.
O desempenho dos estudantes dentro de cada Área/Subárea foi classificado em quartos ( 1/4 ) para os concluintes e para os ingressantes em separado ( ver Figura 1 para o esquema de representação gráfica destas grandezas ).
As notas finais foram colocadas em ordem ascendente e no primeiro quarto foram colocadas as 25 \% notas mais baixas, limitadas superiormente pelo percentil 25 que vem a ser o primeiro quartil.
No segundo quarto, as 25 \% seguintes, limitadas superiormente pelo percentil 50 que vem a ser a mediana.
No terceiro quarto, as 25 \% notas após as encontradas no segundo quarto, ou seja, as primeiras 25 \% acima da mediana, limitadas superiormente pelo percentil 75 que vem a ser o terceiro quartil.
No último quarto foram alocadas as 25 \% maiores notas.
Potencialmente, poderia existir uma relação entre a nota do estudante e a sua percepção sobre a prova.
Calculando as médias das respostas segundo os quartos de notas, seria possível reconhecer a existência ( ou não ) de algum gradiente nas respostas, por exemplo, uma melhor nota correspondendo a uma percepção menos negativa.
Este gradiente seria traduzido, por exemplo, no gráfico 4.1a, pela diminuição das áreas \textbf{laranja} e vermelha ( Difícil e Muito difícil ) com o aumento do desempenho ( indo do primeiro ao último quarto ). 6 Curso, neste relatório, refere -se à unidade de análise para o Conceito ENADE e é caracterizado pela combinação de IES, Área/Subárea ( quando pertinente ) e município de habilitação.
A medida tradicional de “ grau de discriminação ” comparou os resultados dos estudantes do primeiro quarto ( 25 \% piores estudantes ) aos do último quarto ( 25 \% melhores ).
Esta comparação, no gráfico disponibilizado, foi idêntica à comparação da primeira barra ( primeiro quarto ) com a última barra ( último quarto ) de cada grupo ingressante ou concluinte nos gráficos “ segundo quartos de desempenho ” ( gráficos com números ímpares neste capítulo ).
Por outro lado, a disponibilização das informações segundo os quartos de nota permitiu verificar o comportamento de dois outros grupos intermediários que na análise somente dos extremos tem que ser suposto monotônico.
A seguir, serão apresentados os principais resultados relativos aos nove itens avaliados. 4.1 Grau de dificuldade da prova 4.1.1 Formação Geral Os Gráficos 4.1a, 4.1b e 4.1c apresentam a distribuição das respostas à questão “ Qual o grau de dificuldade desta prova na parte de Formação Geral? ”, segundo quartos de nota final e grupo de estudantes concluintes e ingressantes que participaram do ENADE/2009 na área de CIÊNCIAS ECONÔMICAS.
No total de concluintes, 18,4 \% declararam que a prova na parte da Formação Geral foi Difícil/Muito difícil, contra um valor para os ingressantes, 20,7 \%.
Tanto entre os concluintes, como entre os ingressantes, a proporção de indivíduos que classificaram esta parte como Difícil/Muito difícil diminui com o desempenho do estudante, como era esperado.
Por outro lado, a proporção que declarou a parte Fácil/Muito fácil foi menor entre os concluintes ( 25,8 \% ) do que entre os ingressantes ( 26,7 \% ).
Não se observou uma variação monotônica significativa desta proporção com o aumento no desempenho, nem entre os concluintes, nem entre os ingressantes.
Entre os concluintes, 55,9 \% classificaram esta parte com um nível médio de dificuldade.
A proporção entre os ingressantes foi ligeiramente menor, 52,6 \%.
Os Gráficos 4.2a e 4.2b apresentam a distribuição das respostas à questão “ Qual o grau de dificuldade desta prova na parte de Formação Geral? ”, segundo a Região onde o curso está situado para os estudantes concluintes e ingressantes que participaram do ENADE/2009 na área de CIÊNCIAS ECONÔMICAS.
Entre os concluintes e ingressantes, a maior proporção de estudantes que considerou a parte de Formação Geral da prova como de dificuldade média entre todas as Regiões brasileiras, foram os da Região Nordeste com 59,7 \% e 58,8 \%, respectivamente.
As Regiões com as menores proporções desta resposta foram respectivamente para os concluintes e ingressantes, a Sudeste ( 54,5 \% ) e a Centro-Oeste ( 48,1 \% ). 4.1.2 Componente Específico Os Gráficos 4.3a, 4.3b e 4.3c apresentam a distribuição das respostas à questão “ Qual o grau de dificuldade desta prova na parte do Componente Específico? ”, segundo quartos de nota final e grupo de estudantes concluintes e ingressantes que participaram do ENADE/2009 na área de CIÊNCIAS ECONÔMICAS.
Em geral, para as desagregações consideradas, os estudantes classificaram o Componente Específico como mais difícil que a Formação Geral.
No total de concluintes, 38,1 \% declararam que a prova na parte do Componente Específico foi Difícil/Muito difícil, sendo um valor muito maior para os ingressantes, 52,1 \%.
Tanto entre os ingressantes quanto entre os concluintes, a proporção de indivíduos que classificaram esta parte como Muito difícil diminuiu com o desempenho até o terceiro quarto.
A proporção que declara Difícil/Muito difícil entre os concluintes diminuiu com o desempenho, variando de 40,9 \% a 33,7 \%.
A proporção que declarou a parte Fácil/Muito fácil foi maior entre os concluintes ( 8,7 \% ) do que entre os ingressantes ( 6,6 \% ).
Entre os concluintes, 53,2 \% classificaram esta parte com um nível Médio de dificuldade.
A proporção entre os ingressantes foi menor, 41,3 \%.
Entre os concluintes, esta proporção aumentou com o desempenho, indo de 49,0 \% no primeiro quarto de notas a 56,1 \% no último quarto.
Os Gráficos 4.4a e 4.3b apresentam a distribuição das respostas à questão “ Qual o grau de dificuldade desta prova na parte do Componente Específico? ”, segundo a Região onde o curso está situado para os estudantes concluintes e ingressantes que participaram do ENADE/2009 na área de CIÊNCIAS ECONÔMICAS.
Entre os concluintes, a maior proporção de estudantes que considerou a parte do Componente Específico da prova como de dificuldade média entre todas as Regiões brasileiras, foi na Região Sudeste com 57,0 \%, já entre os ingressantes, foram os estudantes da Região Nordeste com 47,1 \%.
A Região com as menores proporções desta resposta foram respectivamente para os concluintes e ingressantes, a Centro-Oeste ( 47,5 \% ) e a Sul ( 37,0 \% ). 4.2 Extensão da prova em relação ao tempo total Os Gráficos 4.5a, 4.5b e 4.5c apresentam a distribuição das respostas à questão sobre a extensão da prova em relação ao tempo total, segundo quartos de nota final e grupo de estudantes concluintes e ingressantes que participaram do ENADE/2009 na área de CIÊNCIAS ECONÔMICAS.
Tanto para concluintes quanto para ingressantes, a maioria considerou a extensão da prova como adequada em relação ao tempo destinado à sua resolução, respectivamente 72,9 \% e 74,2 \%.
Entre os ingressantes, esta proporção cresceu com o desempenho.
A proporção dos que consideraram a extensão da prova muito curta foi idêntica entre os concluintes e ingressantes ( 1,1 \% ).
A proporção dos que consideraram a extensão da prova Curta/Muito curta muito próxima entre os concluintes ( 6,1 \% ) e ingressantes ( 6,0 \% ).
A maior proporção de concluintes e ingressantes declarando que a prova foi curta/muito curta foi entre os estudantes de melhor desempenho, com 9,0 \% e 8,6 \% respectivamente.
A proporção dos que consideraram a extensão da prova muito longa foi idêntica entre os concluintes e ingressantes ( 6,7 \% ).
A proporção dos que consideraram a extensão da prova longa foi maior entre os concluintes ( 14,2 \% ) do que entre os ingressantes ( 13,1 \% ).
Os Gráficos 4.6a e 4.6b apresentam a distribuição das respostas à questão sobre a extensão da prova em relação ao tempo total, segundo a Região onde o curso está situado para os estudantes concluintes e ingressantes que participaram do ENADE/2009 na área de CIÊNCIAS ECONÔMICAS.
Os Gráficos 4.6a e 4.6b apresentam a distribuição das respostas à questão sobre a extensão da prova em relação ao tempo total, segundo a Região onde o curso está situado para os estudantes concluintes e ingressantes que participaram do ENADE/2009 na área de CIÊNCIAS ECONÔMICAS. 4.3 Compreensão dos Enunciados das Questões 4.3.1 Formação Geral Os Gráficos 4.7a, 4.7b e 4.7c apresentam a distribuição das respostas à questão “ Os enunciados das questões da prova na parte de Formação Geral, estavam claros e objetivos? ”, segundo quartos de nota final e grupo de estudantes concluintes e ingressantes que participaram do ENADE/2009 na área de CIÊNCIAS ECONÔMICAS.
Com relação aos enunciados das questões da parte de Formação Geral, as opiniões foram positivas : a grande maioria dos concluintes ( 81,0 \% ) e ingressantes ( 82,6 \% ) consideraram claros e objetivos os enunciados de todas ou da maioria das questões.
Os Gráficos 4.8a e 4.8b apresentam a distribuição das respostas à questão “ Os enunciados das questões da prova na parte de Formação Geral, estavam claros e objetivos? ”, segundo a Região onde o curso está situado para os estudantes concluintes e ingressantes que participaram do ENADE/2009 na área de CIÊNCIAS ECONÔMICAS.
Em termos regionais, as proporções de concluintes que consideraram claros e objetivos os enunciados de todas ou da maioria das questões não variaram de 78,8 \% no Sul a 85,2 \% no Norte.
Já entre os ingressantes estas proporções variaram de 81,9 \% no Sul a 86,2 \% no Centro-Oeste.
Com relação as respostas “ poucos ” e “ não, nenhum ” tanto para concluintes como para ingressantes as maiores proporções ocorreram nas Regiões Sul e Norte, respectivamente 9,7 \% e 7,9 \%. 4.3.2 Componente Específico Os Gráficos 4.9a, 4.9b e 4.9c apresentam a distribuição das respostas à questão “ Os enunciados das questões da prova na parte de formação específica estavam claros e objetivos? ”, segundo quartos de nota final e grupo de estudantes concluintes e ingressantes que participaram do ENADE/2009 na área de CIÊNCIAS ECONÔMICAS.
Com relação aos enunciados das questões da parte do Componente Específico, as opiniões foram também positivas : a grande maioria dos concluintes ( 78,9 \% ) e ingressantes ( 76,1 \% ) consideraram claros e objetivos os enunciados de todas ou da maioria das questões.
A proporção de concluintes declarando que todos os enunciados estavam claros e objetivos apresentou valores para todos os quartos muito próximos ( maior distância 0,8 pontos percentuais ) a da média de 19,7 \% para os concluintes.
Uma situação similar a anterior ocorre para os ingressantes ( maior distância 1,3 pontos percentuais ) com média de 19,5 \%.
Já a combinação de todas ou da maioria das questões para os concluintes foi crescente com o desempenho, indo de 73,0 \% para o primeiro quarto a 83,7 \% para o último quarto.
Entre os ingressantes, esta proporção cresceu com o mesmo comportamento : indo de 70,7 \% no primeiro quarto para 85,1 \% para o grupo com melhor desempenho.
O número de estudantes que declararam que “ não, nenhum ” dos enunciados das questões da parte do Componente Específico era claro e objetivo, foi : 1,1 \% entre os concluintes e 1,7 \% entre os ingressantes.
Os que responderam poucos são 6,1 \% entre os concluintes e 7,5 \% entre os ingressantes.
A proporção de estudantes que respondeu que cerca de metade dos enunciados da parte do Componente Específico era claro e objetivo decresceu com o desempenho tanto para concluintes quanto para ingressantes.
Os Gráficos 4.10a e 4.10b apresentam a distribuição das respostas à questão “ Os enunciados das questões da prova na parte de formação específica estavam claros e objetivos? ”, segundo a Região onde o curso está situado para os estudantes concluintes e ingressantes que participaram do ENADE/2009 na área de CIÊNCIAS ECONÔMICAS.
Em termos regionais, as proporções de ingressantes e de concluintes que consideraram claros e objetivos os enunciados de todas ou da maioria das questões variaram mais entre os ingressantes ( no máximo 5,6 pontos percentuais ) do que entre os concluintes ( no máximo 2,9 pontos percentuais ).
Entre os concluintes, os valores foram de 78,0 \% no Centro-Oeste a 80,9 \% no Nordeste.
Já entre os ingressantes estas proporções variaram de 73,7 \% na Região Norte a 79,3 \% no Centro-Oeste.
A Região Sudeste, também, apresentou valores ( 76,5 \% ) acima da média ( 76,0 \% ).
Para os ingressantes, as Regiões que apresentaram valores acima da média ( 9,1 \% ) foram a Norte, Nordeste e a Sul, com valores respectivamente de 10,0 \%, 9,3 \% e 9,8 \%. 4.4 Suficiência das informações/instruções fornecidas Os Gráficos 4.11a, 4.11b e 4.11c apresentam a distribuição das respostas à questão “ As informações/instruções fornecidas para a resolução das questões foram suficientes para resolvê -las? ”, segundo quartos de nota final e grupo de estudantes concluintes e ingressantes que participaram do ENADE/2009 na área de CIÊNCIAS ECONÔMICAS.
Quando perguntados sobre as informações/instruções fornecidas para a resolução das questões, quase 86,9 \% dos concluintes e 82,8 \% dos ingressantes de todo o Brasil responderam que estas eram ou excessivas ou suficientes em todas ou na maioria das questões.
Tal proporção foi crescente com o desempenho, tanto para concluintes como para ingressantes.
Entre os concluintes a proporção vai de 82,1 \% para o primeiro quarto a 90,6 \% para o último quarto.
Entre os ingressantes as proporções equivalentes foram, respectivamente, 77,6 \% a 91,0 \%.
O complemento desta opinião, ou seja, os estudantes que responderam somente em algumas ou em nenhuma delas, apresentou uma proporção decrescente com o desempenho tanto para os concluintes como para os ingressantes.
Os Gráficos 4.12a e 4.12b apresentam a distribuição das respostas à questão “ As informações/instruções fornecidas para a resolução das questões foram suficientes para resolvê -las? ”, segundo a Região onde o curso está situado para os estudantes concluintes e ingressantes que participaram do ENADE/2009 na área de CIÊNCIAS ECONÔMICAS.
Os Gráficos 4.12a e 4.12b apresentam a distribuição das respostas à questão “ As informações/instruções fornecidas para a resolução das questões foram suficientes para resolvê -las? ”, segundo a Região onde o curso está situado para os estudantes concluintes e ingressantes que participaram do ENADE/2009 na área de CIÊNCIAS ECONÔMICAS. 4.5 Dificuldade encontrada ao responder à prova Os Gráficos 4.13a, 4.13b e 4.13c apresentam a distribuição das respostas à questão “ Você se deparou com alguma dificuldade ao responder a prova? ”, segundo quartos de nota final e grupo de estudantes concluintes e ingressantes que participaram do ENADE/2009 na área de CIÊNCIAS ECONÔMICAS.
Perguntados sobre as dificuldades com as quais se depararam ao responder à prova, 59,3 \% dos ingressantes apontaram como principal dificuldade o desconhecimento do conteúdo, já os 33,4 \% dos concluintes apontaram como principal dificuldade a forma diferente de abordagem.
O desconhecimento do conteúdo foi apontado somente por 17,3 \% dos concluintes.
A forma diferente de abordagem do conteúdo foi considerada como maior dificuldade por 16,1 \% dos ingressantes.
Já a falta de motivação foi a principal dificuldade para 29,9 \% dos concluintes e para 15,5 \% dos ingressantes.
Esta dificuldade apresentou valores crescentes conforme o desempenho para os concluintes.
Em termos nacionais, uma proporção maior de concluintes ( 13,0 \% ) do que de ingressantes ( 5,4 \% ) declarou que não teve qualquer dificuldade para responder à prova.
Os Gráficos 4.14a e 4.14b apresentam a distribuição das respostas à questão “ Você se deparou com alguma dificuldade ao responder a prova? ”, segundo a Região onde o curso está situado para os estudantes concluintes e ingressantes que participaram do ENADE/2009 na área de CIÊNCIAS ECONÔMICAS.
Em relação às Regiões brasileiras, o percentual de ingressantes que apontaram o desconhecimento do conteúdo como principal dificuldade ao responder à prova variou de 55,9 \% na Região Nordeste, a 60,6 \% na Centro-Oeste.
Entre os concluintes, tal percentual variou de 14,3 \% na Região Sudeste a 23,4 \%, na Centro-Oeste.
A forma diferente de abordagem de conteúdo foi citada como maior dificuldade por 28,9 \% dos concluintes na Região Sudeste chegando a 41,3 \% na Região Norte.
Já entre os ingressantes os valores foram de 14,1 \% na Região Sudeste a 21,8 \% na Região Norte.
O percentual de concluintes que citaram a falta de motivação como principal dificuldade variou de 22,0 \% na Região Norte a 34,5 \% na Região Sudeste.
Entre os ingressantes as estatísticas correspondentes foram de 11,3 \% na Região Norte a 18,1 \%, na Região Nordeste.
Os que declararam que não tiveram qualquer tipo de dificuldade variaram, entre os concluintes, de 10,1 \% na Região Centro-Oeste a 14,8 \% na Região Sudeste.
Entre os ingressantes este percentual variou de 3,8 \% na Região Centro-Oeste a 6,4 \% na Região Sul. 4.6 Percepção sobre os conteúdos das questões objetivas da prova Os Gráficos 4.15a, 4.15b e 4.15c apresentam a distribuição das respostas à questão “ Considerando apenas as questões objetivas da prova, você percebeu que :... ”, segundo quartos de nota final e grupo de estudantes concluintes e ingressantes que participaram do ENADE/2009 na área de CIÊNCIAS ECONÔMICAS.
As respostas possíveis para esta questão estavam relacionadas à cobertura do conteúdo das questões objetivas durante o curso.
Como esperado, a proporção de ingressantes ( 65,5 \% ) que declarou que não estudou ainda a maioria desses conteúdos foi bem maior do que entre os concluintes ( 3,1 \% ).
Esta proporção foi decrescente para os concluintes como função do desempenho.
Por outro lado, a proporção de estudantes que declarou que estudou e aprendeu todos ou a maioria desses conteúdos, entre os ingressantes foi de 16,5 \% e entre os concluintes 60,5 \%.
Esta proporção foi crescente como função do desempenho, tanto para os concluintes como para os ingressantes.
A proporção de estudantes que declarou ter estudado a maioria desses conteúdos, mas não os ter aprendido foi bem maior entre os concluintes ( 24,9 \% ) do que os ingressantes ( 5,5 \% ).
Os Gráficos 4.16a e 4.16b apresentam a distribuição das respostas à questão “ Considerando apenas as questões objetivas da prova, você percebeu que :... ”, segundo a Região onde o curso está situado para os estudantes concluintes e ingressantes que participaram do ENADE/2009 na área de CIÊNCIAS ECONÔMICAS.
Entre os ingressantes, os percentuais variaram de 63,2 \% na Região Sudeste a 69,8 \%, na Norte.
Em todas as Regiões, a maioria dos concluintes afirmou ter estudado e aprendido todos ou muitos dos conteúdos, com percentuais que variaram de 50,2 \% na Região Centro-Oeste a 67,0 \%, na Sudeste.
Para os ingressantes, tal percentual variou de 13,8 \%, nas Regiões Norte e Sul, a 19,1 \%, na Sudeste. 4.7 Tempo gasto para concluir a prova Os Gráficos 4.17a, 4.17b e 4.17c apresentam a distribuição das respostas à questão “ Qual foi o tempo gasto para concluir a prova? ”, segundo quartos de nota final e grupo de estudantes concluintes e ingressantes que participaram do ENADE/2009 na área de CIÊNCIAS ECONÔMICAS.
Quando perguntados quanto ao tempo de conclusão da prova, a maioria respondeu entre uma e duas horas, tanto entre os concluintes ( 44,5 \% ) quanto entre os ingressantes ( 48,9 \% ).
A segunda maior incidência de respostas foi entre duas e três horas : concluintes ( 34,0 \% ) e ingressantes ( 29,8 \% ).
A proporção daqueles que declararam ter usado as quatro horas mas não ter conseguido terminar foi bem baixo e idêntica a 1,1 \% tanto entre os concluintes quanto entre os ingressantes.
Os Gráficos 4.18a e 4.18b apresentam a distribuição das respostas à questão “ Qual foi o tempo gasto por você para concluir a prova? ”, segundo a Região onde o curso está situado para os estudantes concluintes e ingressantes que participaram do ENADE/2009 na área de CIÊNCIAS ECONÔMICAS.
A proporção de concluintes que gastaram entre uma e duas horas foi maior na Região Sudeste ( 46,7 \% ) e menor na Região Centro-Oeste ( 35,7 \% ).
Para ingressantes, o maior percentual para este tempo ocorreu na Região Sul ( 52,6 \% ) e a menor na Região Norte ( 45,5 \% ).

% matched lemmas: laranja
\end{textsample}
